% Generated draft: requires mathematical and source review.
\documentclass[11pt]{article}
\usepackage[margin=25mm]{geometry}
\usepackage{fontspec}
\setmainfont{Latin Modern Roman}
\usepackage{amsmath,amssymb,mathtools,unicode-math}
\setmathfont{Latin Modern Math}
\usepackage{enumitem,needspace,xurl,hyperref}
\setlength{\parindent}{0pt}
\setlength{\parskip}{5pt}
\setlength{\emergencystretch}{3em}
\newcommand{\N}{\mathbb{N}}
\newcommand{\Z}{\mathbb{Z}}
\newcommand{\Q}{\mathbb{Q}}
\newcommand{\R}{\mathbb{R}}
\newcommand{\C}{\mathbb{C}}
\newcommand{\F}{\mathbb{F}}
\newcommand{\A}{\mathbb{A}}
\newcommand{\PP}{\mathbb{P}}
\newcommand{\E}{\mathbb{E}}
\newcommand{\eps}{\varepsilon}
\newcommand{\dd}{\,\mathrm{d}}
\begin{document}
\section*{UnsolvedMath ID: 30000592}
\textbf{Source problem code:} catalogue entry
\subsection{Definitions and mathematical statement}
Let \(M\) be a closed smooth spin manifold of dimension \(m\), with its spin structure fixed, and put \(\Gamma=\pi_1(M)\). In the real \(K\)-theory convention required by the displayed \(KO\)-groups, use the real maximal and reduced group \(C^*\)-algebras. The Rosenberg indices are the indices of the spin Dirac operator twisted by the corresponding flat Mishchenko bundles:
\[
\alpha_{\max}(M)\in KO_m(C^*_{\max,\R}\Gamma),\qquad
\alpha_{\mathrm{red}}(M)\in KO_m(C^*_{\mathrm{red},\R}\Gamma).
\]
The degree is taken modulo eight. The quotient homomorphism \(q:C^*_{\max,\R}\Gamma\to C^*_{\mathrm{red},\R}\Gamma\) gives the comparison
\[
q_*\alpha_{\max}(M)=\alpha_{\mathrm{red}}(M).
\]
Thus nonvanishing of the reduced index implies nonvanishing of the maximal one; the converse requires information about this map and these classes.

A positive-scalar-curvature metric means a smooth Riemannian metric with scalar curvature strictly positive everywhere. The standard index obstruction gives the necessary condition
\[
M\text{ admits such a metric}\ \Longrightarrow\
\alpha_{\max}(M)=0\quad\text{and}\quad\alpha_{\mathrm{red}}(M)=0.
\]
The source asks how much of the obstruction theory is captured by these indices and how the two versions compare. It does not give a precise general sufficiency statement. One must not replace that broad question by an asserted equivalence between vanishing and existence for all closed spin manifolds.

A definite comparison theorem would need to specify the dimension, fundamental-group assumptions, stable versus direct existence of a metric on \(M\), and which additional obstructions or sufficiency conditions are to be tested. No surgery-range dimension assumption, torsion-free restriction, or stabilization is imposed from context. The extract also suppresses the real-algebra and connectedness conventions; for disconnected manifolds the formulation is ordinarily interpreted componentwise.

The supplied assessment checked 2026-08-21 reports important special-case index-theoretic results but no complete answer to the broad information question. It does not reproduce the hypotheses or conclusions of the cited enlargeability work, which are therefore not promoted to a general characterization.
\subsection{Short English statement}
Determine which positive-scalar-curvature obstructions the maximal and reduced Rosenberg indices detect, and what their vanishing can imply under precise hypotheses. The source asks a broad comparison, not a universally asserted vanishing criterion for existence.
\subsection{Sources}
\begin{itemize}
\item Local UnsolvedMath registry, permanent ID 30000592, catalogue entry, \emph{Positive Scalar Curvature Information from Rosenberg Indices}.
\item \emph{Topologie} (2007), Oberwolfach Reports, pp. 2579--2630, DOI \url{https://doi.org/10.4171/owr/2006/43}.
\item Hanke--Schick, supplied label \emph{Enlargeability and index theory}: \url{https://www.uni-math.gwdg.de/schick/publ/Enlargeability_and_index_theory_final_for_JDG.pdf}.
\item Supplied literature triage checked 2026-08-21; no complete general sufficiency criterion or exact special-case hypotheses are reproduced.
\end{itemize}
\end{document}
